% Заголовок, автор, дата
\title{Экспертное исследование}
\author{Ивасюк С.И. \\ Межрегиональная общественная организация "Национальный союз землепользователей"}
\date{10 апреля 2025 г.}
\maketitle

%  Аннотация
\begin{abstract}
Экспертное исследование направлено на выявление возможностей оформления пользования водными объектами, расположенными на землях сельскохозяйственного назначения. Особое внимание уделено  
\end{abstract}

% Содержание
\tableofcontents


\section{Введение}
\label{sec:intro}

В этом разделе представлено краткое введение в проблематику исследования. Современные научные публикации требуют четкой структуры и профессионального оформления.

\section{Основная часть}
\label{sec:main}

\subsection{Пример изображения}
\label{subsec:image}

\begin{figure}[H]
    \centering
    \includegraphics[width=0.5\textwidth]{example-image}
    \caption{Пример эталонного изображения}
    \label{fig:example-ref}
\end{figure}

% Исправленные вставки (теперь с правильными подписями):
% Стандартное использование
\insertimage{example-image.pdf}{Схема работы алгоритма}

% С дополнительными параметрами
\insertimage[width=0.5\textwidth,angle=15]{photo.jpg}{Пример повернутого изображения}

% Безопасная версия (псевдоним)
\safeinsertimage[width=0.8\textwidth]{diagram.jpg}{Архитектура системы}

На рис.~\ref{fig:example-ref} показан базовый пример.

\subsection{Примеры списков}
\label{subsec:lists}

\subsubsection{Неупорядоченный список}
\begin{itemize}
    \item Первый элемент списка
    \item Второй элемент списка
    \item Третий элемент списка
\end{itemize}

\subsubsection{Многоуровневый список без нумерации и маркировки}
\begin{description}
    \item[Первый уровень] 
        \begin{description}
            \item[Второй уровень] 
                \begin{description}
                    \item[Третий уровень] Пример текста
                \end{description}
        \end{description}
    \item[Другой элемент] Текст описания
\end{description}

\paragraph{Параграф} Пример параграфа в документе. Текст параграфа может содержать различные элементы оформления.

\subparagraph{Подпараграф} Пример подпараграфа с дополнительными уточнениями.

\subsection{Математические выражения}
\label{subsec:math}

Пример математического выражения в тексте: $E = mc^2$.

Уравнение в выключной формуле:
\begin{equation}
    \label{eq:1}
    \int_{-\infty}^\infty e^{-x^2} dx = \sqrt{\pi}
\end{equation}

Система уравнений:
\[
\begin{cases}
    \frac{dx}{dt} = \sigma(y - x) \\
    \frac{dy}{dt} = x(\rho - z) - y \\
    \frac{dz}{dt} = xy - \beta z
\end{cases}
\]

\subsection{Пример таблицы}
\label{subsec:table}

\begin{table}[h]
    \centering
    \caption{Пример таблицы с границами}
    \label{tab:example}
    \begin{tabular}{|l|c|r|}
        \hline
        Заголовок 1 & Заголовок 2 & Заголовок 3 \\
        \hline
        Данные 1 & Данные 2 & Данные 3 \\
        \hline
        Данные 4 & Данные 5 & Данные 6 \\
        \hline
    \end{tabular}
\end{table}

В таблице \ref{tab:example} представлен пример оформления таблицы с границами.

\section{Заключение}
\label{sec:conclusion}

В данной статье были рассмотрены основные элементы оформления научной статьи с использованием системы \LaTeX. Приведенные примеры могут служить шаблоном для подготовки собственных публикаций.
